\documentclass[12pt,a4paper]{article}

% ── Packages ──────────────────────────────────────────────
\usepackage[utf8]{inputenc}
\usepackage[T1]{fontenc}
\usepackage{lmodern}
\usepackage{amsmath,amssymb,amsthm}
\usepackage{mathtools}
\usepackage{geometry}
\geometry{margin=1in}
\usepackage{graphicx}
\graphicspath{{figures/}}
\usepackage{booktabs}
\usepackage{hyperref}
\hypersetup{
    colorlinks=true,
    linkcolor=blue!60!black,
    citecolor=blue!60!black,
    urlcolor=blue!60!black
}
\usepackage{caption}
\usepackage{float}
\usepackage{enumitem}
\usepackage{xcolor}

% ── Custom commands ───────────────────────────────────────
\newcommand{\Hfour}{H_4}
\newcommand{\Hthree}{H_3}
\newcommand{\Htwo}{H_2}
\newcommand{\Eeight}{E_8}
\newcommand{\vphi}{\varphi}
\newcommand{\Lmat}{L_{H_4}}
\newcommand{\R}{\mathbb{R}}
\newcommand{\Q}{\mathbb{Q}}
\newcommand{\Z}{\mathbb{Z}}

\newtheorem{definition}{Definition}
\newtheorem{proposition}{Proposition}
\newtheorem{observation}{Observation}

% ── Title ─────────────────────────────────────────────────
\title{%
  \textbf{Nonlinear Mode Localisation on the $\Hfour$ Graph}\\[6pt]
  \large Part~III of a Series on Symmetry-Constrained Resonant Structure%
}
\author{%
  Lee Smart\\[4pt]
  \textit{Institute of Vibrational Field Dynamics}\\[2pt]
  \texttt{contact@vibrationalfielddynamics.org}\\[2pt]
  \texttt{@vfd\_org}%
}
\date{\today}

% ══════════════════════════════════════════════════════════
\begin{document}
\maketitle
\thispagestyle{empty}

% ── Abstract ──────────────────────────────────────────────
\begin{abstract}
The nonlinear dynamics of a real scalar field on the 600-cell vertex
graph are studied under variation of the self-interaction strength~$\beta$.
The system admits a Hamiltonian formulation with an $\Hfour$-symmetric
graph Laplacian coupling and a quartic on-site nonlinearity.
Part~I of this series established that the system exhibits persistent
spatial localisation and coherent shell dynamics; Part~II identified
seven structural invariants of the Laplacian spectrum that constrain
the available dynamical channels.
The present work classifies the space of persistent nonlinear
configurations.  Four attractor classes are identified: backbone
harmonic modes, phase-locked multi-mode states, spatially localised
breather-like configurations, and transitional unstable states.
These classes are characterised by their inverse participation ratio,
spectral composition, recurrence statistics, and stability under
perturbation.  A phase diagram in~$\beta$ delineates the boundaries
between linear, weakly nonlinear, and strongly nonlinear regimes.
The attractor taxonomy is shown to be consistent with the spectral
invariants of Part~II: the seven structural constraints restrict the
attractor space to a finite, classifiable set of dynamical families.
All results are derived from numerical simulation with no physical
interpretation assumed.
\end{abstract}

\vspace{1em}
\noindent\textbf{Keywords:}
nonlinear lattice dynamics, discrete breathers, attractor classification,
$\Hfour$ symmetry, 600-cell, spectral constraints, inverse participation
ratio, Hamiltonian systems.

\newpage
\tableofcontents
\newpage


% ══════════════════════════════════════════════════════════
\section{Introduction}\label{sec:intro}
% ══════════════════════════════════════════════════════════

In Part~I of this series~\cite{paper1}, a real scalar field was defined
on the 120-vertex graph of the 600-cell and evolved under a nonlinear
wave equation incorporating the $\Hfour$ graph Laplacian.  Numerical
simulations revealed persistent spatial localisation, coherent shell
reflection, and recurrent breathing dynamics that are absent in
degree-matched control graphs.  These phenomena were attributed to the
extreme spectral compression of the $\Hfour$ Laplacian, which produces
only nine distinct eigenvalues across 120~degrees of freedom.

In Part~II~\cite{paper2}, the internal organisation of this nine-sector
spectrum was analysed algebraically.  Seven structural invariants were
identified---$\vphi$-cancellation, degree anchoring, conjugate pairing,
square multiplicity folding, three-zone partition, spectral backbone
decomposition, and departure structure---collectively demonstrating
that the spectrum is governed by a highly constrained algebraic
organisation rather than arising generically.

The present work addresses a natural question arising from these
results: \emph{given a spectrum with only nine sectors and seven
governing invariants, what is the space of stable nonlinear dynamical
configurations?}

This is a question about dynamical attractors, not about physical
entities.  The goal is to classify the persistent configurations of the
nonlinear field system---that is, the states to which generic initial
conditions evolve and in which the system remains for long times.  The
classification is based on observable dynamical properties: spatial
localisation, spectral composition, temporal periodicity, and stability
under perturbation.

The distinction between spectral structure and dynamical structure is
central.  The spectral invariants of Part~II are properties of the
static Laplacian; they constrain the \emph{channels} available to
dynamics but do not determine a unique dynamical outcome.  The
attractor classification of the present work identifies which of these
constrained channels are actually realised as persistent nonlinear
configurations.

\paragraph{Organisation.}
Section~\ref{sec:system} defines the full nonlinear system and its
conserved quantities.  Section~\ref{sec:taxonomy} introduces the mode
taxonomy: single-mode states, multi-mode locked states, and localised
breather-like states.  Section~\ref{sec:stability} presents stability
analysis including persistence, perturbation response, and parameter
sweeps.  Section~\ref{sec:quantify} defines quantitative diagnostics.
Section~\ref{sec:invariants} maps the attractor classes to the
spectral invariants of Part~II.  Section~\ref{sec:classification}
presents the complete attractor classification.
Section~\ref{sec:numerics} reports the numerical results.
Section~\ref{sec:discussion} discusses the results in relation to
spectral structure, localisation mechanisms, and limitations.
Section~\ref{sec:conclusion} offers conclusions.


% ══════════════════════════════════════════════════════════
\section{Dynamical System Setup}\label{sec:system}
% ══════════════════════════════════════════════════════════

\subsection{Equation of Motion}

The system studied throughout this series is defined by the nonlinear
wave equation on the 600-cell vertex graph:
\begin{equation}\label{eq:eom}
  \ddot{\Phi}_i + \omega_0^2 \, \Phi_i
  + \lambda \, (\Lmat \Phi)_i
  + \beta \, \Phi_i^3 = 0,
  \qquad i = 1, \ldots, 120,
\end{equation}
where $\Phi_i(t) \in \R$ is the field amplitude at vertex~$i$,
$\Lmat = 12I - A$ is the graph Laplacian of the 600-cell,
$\omega_0$ is the base frequency, $\lambda$ is the Laplacian coupling
strength, and $\beta$ is the nonlinear self-interaction parameter.

Parts~I and~II studied this system primarily in the linear regime
($\beta = 0$) and with moderate nonlinearity.  The present work
treats $\beta$ as the primary control parameter and maps the full
range from linear ($\beta \to 0$) through strongly nonlinear
behaviour.

\subsection{Hamiltonian Structure}

Equation~\eqref{eq:eom} derives from the Hamiltonian
\begin{equation}\label{eq:ham}
  H = \sum_{i=1}^{120} \left[
    \frac{1}{2} \Pi_i^2
    + \frac{1}{2} \omega_0^2 \, \Phi_i^2
    + \frac{1}{4} \beta \, \Phi_i^4
  \right]
  + \frac{1}{2} \lambda \, \Phi^T \Lmat \Phi,
\end{equation}
where $\Pi_i = \dot{\Phi}_i$ is the conjugate momentum.  The
equations of motion are
\begin{equation}\label{eq:hamilton_eqs}
  \dot{\Phi}_i = \Pi_i, \qquad
  \dot{\Pi}_i = -\omega_0^2 \, \Phi_i
    - \lambda \, (\Lmat \Phi)_i - \beta \, \Phi_i^3.
\end{equation}
The Hamiltonian~$H$ is a conserved quantity: $\dot{H} = 0$.  This
conservation provides a fundamental check on the numerical integration.

\subsection{Phase Space}

The phase space of the system is $(\Phi, \Pi) \in \R^{240}$, where
$\Phi = (\Phi_1, \ldots, \Phi_{120})$ and
$\Pi = (\Pi_1, \ldots, \Pi_{120})$.

The $\Hfour$ symmetry group acts on phase space by simultaneous
permutation of vertex indices in both~$\Phi$ and~$\Pi$.  An orbit of
this action corresponds to a set of dynamically equivalent
configurations related by the symmetry of the 600-cell.

\subsection{Observables}\label{sec:observables}

The following observables are used throughout this work to characterise
dynamical states:

\begin{enumerate}[nosep]
\item \textbf{Total energy:}
  $H$ as defined in~\eqref{eq:ham}.

\item \textbf{Per-vertex energy:}
  $h_i = \frac{1}{2}\Pi_i^2 + \frac{1}{2}\omega_0^2 \Phi_i^2
  + \frac{1}{4}\beta \Phi_i^4
  + \frac{1}{2}\lambda \Phi_i (\Lmat \Phi)_i$.

\item \textbf{Mode projection amplitudes:}
  $c_k(t) = \langle \Phi(t), \mathbf{v}^{(k)} \rangle$, where
  $\mathbf{v}^{(k)}$ are the eigenvectors of~$\Lmat$.

\item \textbf{Spatial variance:}
  $\sigma^2(t) = \frac{1}{120} \sum_i (\Phi_i - \bar{\Phi})^2$,
  where $\bar{\Phi} = \frac{1}{120} \sum_i \Phi_i$.
\end{enumerate}


% ══════════════════════════════════════════════════════════
\section{Mode Taxonomy}\label{sec:taxonomy}
% ══════════════════════════════════════════════════════════

The nonlinear system~\eqref{eq:eom} admits several qualitatively
distinct classes of persistent dynamical configuration.  These are
distinguished by their spectral content, spatial profile, and temporal
behaviour.

% ── 3.1 ──────────────────────────────────────────────────
\subsection{Single-Mode States}\label{sec:single}

\begin{definition}[Single-mode state]
A single-mode state is a solution dominated by projection onto a
single eigenvector of the Laplacian:
\begin{equation}\label{eq:single}
  \Phi(t) \approx A_k \cos(\Omega_k t) \, \mathbf{v}^{(k)},
\end{equation}
where $\mathbf{v}^{(k)}$ is an eigenvector of~$\Lmat$ with eigenvalue
$\mu_k$, and $\Omega_k$ is the effective nonlinear frequency.
\end{definition}

In the linear limit ($\beta = 0$), the frequency is
$\Omega_k^2 = \omega_0^2 + \lambda \mu_k$.  The nonlinear correction
shifts this to
\begin{equation}\label{eq:freq_shift_single}
  \Omega_k^2 = \omega_0^2 + \lambda \mu_k
  + \frac{3}{4} \beta A_k^2 \, \lVert \mathbf{v}^{(k)} \rVert_4^4,
\end{equation}
where $\lVert \mathbf{v}^{(k)} \rVert_4^4
= \sum_i |v_i^{(k)}|^4$ is the $\ell^4$ norm of the eigenvector.
The factor~$3/4$ arises from time-averaging the cubic term over one
period.

Single-mode states inherit the spatial symmetry of their constituent
eigenvector.  Since the eigenspaces of~$\Lmat$ decompose into
irreducible representations of~$\Hfour$, single-mode states carry
definite transformation properties under the symmetry group.

\begin{observation}
Single-mode states associated with backbone sectors
($k = 1, \ldots, 6$, comprising 91~modes) are more stable than those
in the departure sectors ($k = 7, 8, 9$), as measured by persistence
time under small perturbation.  This is consistent with the spectral
backbone hierarchy identified in Part~II (Invariant~6).
\end{observation}

% ── 3.2 ──────────────────────────────────────────────────
\subsection{Multi-Mode Locked States}\label{sec:multi}

\begin{definition}[Multi-mode locked state]
A multi-mode locked state is a superposition of eigenmodes that
maintains fixed phase relationships over long times:
\begin{equation}\label{eq:multi}
  \Phi(t) = \sum_{k \in \mathcal{K}} c_k(t) \, \mathbf{v}^{(k)},
\end{equation}
where $\mathcal{K} \subset \{1, \ldots, 120\}$ is a subset of mode
indices and the relative phases $\theta_{kl}(t) = \arg(c_k) - \arg(c_l)$
are approximately constant: $|\dot{\theta}_{kl}| < \epsilon$ for all
$k, l \in \mathcal{K}$.
\end{definition}

The nonlinearity~$\beta \Phi^3$ couples different eigenmodes.  For a
multi-mode superposition, the effective nonlinear frequency of mode~$k$
becomes
\begin{equation}\label{eq:freq_shift_multi}
  \Omega_k^2 \to \omega_0^2 + \lambda \mu_k
  + \frac{3}{2} \beta |c_k|^2 \, \gamma_{kk}
  + 3\beta \sum_{l \neq k} |c_l|^2 \, \gamma_{kl},
\end{equation}
where $\gamma_{kl} = \sum_i v_i^{(k)} v_i^{(k)} v_i^{(l)} v_i^{(l)}$
are overlap integrals determined by the spatial profiles of the
eigenvectors.  Phase locking occurs when the nonlinear frequency shifts
bring modes into a rational frequency ratio.

\begin{observation}
Stable multi-mode locked states preferentially involve modes from the
same $\vphi$-zone (Invariant~5 of Part~II): either both from
$\mathcal{S}_-$, both from $\mathcal{S}_0$, or both from
$\mathcal{S}_+$.  Cross-zone locking is observed but is less stable,
consistent with the three-zone partition constraining the coupling
structure.
\end{observation}

The resonance locking condition for two modes $k$ and $l$ is
\begin{equation}\label{eq:lock}
  \frac{\Omega_k}{\Omega_l} = \frac{p}{q}, \qquad p, q \in \Z_{>0},
\end{equation}
where $p/q$ is a low-order rational.  The $\Hfour$ spectral compression
limits the number of available frequency ratios to those derivable from
nine sectors, producing a sparse resonance network compared to generic
graphs.

% ── 3.3 ──────────────────────────────────────────────────
\subsection{Localised (Breather-Like) States}\label{sec:breather}

\begin{definition}[Breather-like state]
A breather-like state is a spatially localised, time-periodic (or
quasi-periodic) configuration satisfying:
\begin{enumerate}[nosep]
  \item High spatial concentration:
    $\mathrm{IPR}(\Phi(t)) \gg 1/120$ for all~$t$.
  \item Energy trapping: a fraction $f > 0.5$ of the total energy
    remains within a connected subgraph of diameter~$\leq 2$ for
    times $t > T_{\mathrm{transient}}$.
  \item Persistence: $\lVert \Phi(t+T) - \Phi(t) \rVert / \lVert
    \Phi(t) \rVert < \delta$ for some period~$T$ and threshold
    $\delta \ll 1$.
\end{enumerate}
\end{definition}

Breather-like states emerge from the interplay between nonlinear
self-trapping (the $\beta \Phi^3$ term) and the discrete geometry of
the 600-cell graph.  In the theory of discrete breathers on
lattices~\cite{flach2008,aubry1997}, spatially localised time-periodic
solutions exist generically in nonlinear Hamiltonian lattice systems
above a nonlinearity threshold.

The 600-cell graph presents a distinctive setting for discrete breather
formation.  The 12-regular structure with icosahedral local geometry
provides a homogeneous coupling environment, while the extreme
spectral compression limits the radiation channels available for
energy dispersion.  The backbone modes (91~of 120~degrees of freedom)
occupy a narrow spectral band, reducing the density of states into
which a localised excitation can radiate.

\begin{observation}
Breather-like states on the 600-cell graph are characterised by:
\begin{itemize}[nosep]
  \item high amplitude concentration on $\mathcal{O}(10)$ vertices,
  \item energy trapping persisting over $> 10^3$ oscillation periods,
  \item broad spectral composition spanning multiple sectors.
\end{itemize}
These properties are consistent with the localisation enhancement
reported in Part~I ($5\times$ enhancement relative to controls).
\end{observation}


% ══════════════════════════════════════════════════════════
\section{Stability Analysis}\label{sec:stability}
% ══════════════════════════════════════════════════════════

The classification of dynamical states requires distinguishing
persistent attractors from transient configurations.  Three
complementary stability measures are employed.

% ── 4.1 ──────────────────────────────────────────────────
\subsection{Persistence Over Time}\label{sec:persist}

For a candidate periodic or quasi-periodic state with approximate
period~$T$, the recurrence error is defined as
\begin{equation}\label{eq:recurrence}
  R(T) = \frac{\lVert \Phi(t_0 + T) - \Phi(t_0) \rVert}
              {\lVert \Phi(t_0) \rVert},
\end{equation}
evaluated after a transient time $t_0 \gg T$.  A state is classified
as \emph{persistent} if $R(T) < \delta_R$ for threshold $\delta_R = 0.05$
over a window of at least 100~periods.

States are further classified by the nature of their recurrence:
\begin{itemize}[nosep]
  \item \textbf{Periodic:} $R(T) < 10^{-3}$ for a single dominant
    period~$T$.
  \item \textbf{Quasi-periodic:} $R(T) < \delta_R$ for a period~$T$
    that is a combination of two or more incommensurate
    frequencies; the power spectrum shows discrete peaks.
  \item \textbf{Weakly chaotic:} $R(T) > \delta_R$ for all~$T$
    within the simulation window, but the largest Lyapunov exponent
    is small ($\lambda_{\max} < 10^{-2}$).
  \item \textbf{Transient:} the state decays or disperses within the
    simulation window.
\end{itemize}

% ── 4.2 ──────────────────────────────────────────────────
\subsection{Robustness to Perturbation}\label{sec:robust}

Given a state $(\Phi^*, \Pi^*)$ identified as persistent, a perturbed
initial condition is constructed:
\begin{equation}\label{eq:perturb}
  \Phi_{\epsilon} = \Phi^* + \epsilon \, \boldsymbol{\eta},
  \qquad
  \Pi_{\epsilon} = \Pi^*,
\end{equation}
where $\boldsymbol{\eta}$ is drawn uniformly from the unit sphere in
$\R^{120}$ and $\epsilon$ controls the perturbation amplitude.

The divergence rate is measured by
\begin{equation}\label{eq:diverge}
  D(t) = \frac{\lVert \Phi_{\epsilon}(t)
    - \Phi_{\mathrm{ref}}(t) \rVert}
    {\lVert \Phi_{\mathrm{ref}}(t) \rVert},
\end{equation}
where $\Phi_{\mathrm{ref}}(t)$ is the unperturbed evolution.

A state is classified as \emph{stable} if $D(t)$ remains bounded for
$t \gg T$, and the system returns to a neighbourhood of the original
attractor.  The basin of attraction is estimated by varying~$\epsilon$
and identifying the critical value $\epsilon_c$ at which the perturbed
trajectory no longer returns.

% ── 4.3 ──────────────────────────────────────────────────
\subsection{Parameter Sweep in $\beta$}\label{sec:sweep}

The nonlinear coupling~$\beta$ is swept across the range
$\beta \in [0, 5]$ with fixed $\omega_0 = 1.0$ and $\lambda = 1.0$.
For each value of~$\beta$, an ensemble of initial conditions is
evolved and classified by the diagnostics of
Section~\ref{sec:quantify}.

Three dynamical regimes are identified:

\begin{enumerate}[nosep]
\item \textbf{Linear regime} ($\beta \lesssim 0.01$):
  dynamics are well-approximated by superpositions of linear
  eigenmodes.  No mode locking or localisation beyond the linear
  spectral structure is observed.

\item \textbf{Weakly nonlinear regime}
  ($0.01 \lesssim \beta \lesssim 1.0$):
  nonlinear frequency shifts produce mode locking.
  Multi-mode locked states appear.  Localisation is enhanced
  relative to the linear case but breather-like states are
  not yet dominant.

\item \textbf{Strongly nonlinear regime} ($\beta \gtrsim 1.0$):
  self-trapping dominates.  Spatially localised breather-like
  states form spontaneously from generic initial conditions.
  Energy concentration on $\mathcal{O}(10)$ vertices is
  typical.
\end{enumerate}

The boundaries between regimes are not sharp; the classification
represents the dominant behaviour within each range.  The phase diagram
(Figure~\ref{fig:phase}) maps the attractor class frequency as a
function of~$\beta$.


% ══════════════════════════════════════════════════════════
\section{Quantifying Structure}\label{sec:quantify}
% ══════════════════════════════════════════════════════════

% ── 5.1 ──────────────────────────────────────────────────
\subsection{Inverse Participation Ratio}\label{sec:ipr}

The inverse participation ratio (IPR) quantifies spatial localisation:
\begin{equation}\label{eq:ipr}
  \mathrm{IPR}(t) = \frac{\sum_{i=1}^{120} \Phi_i(t)^4}
    {\left(\sum_{i=1}^{120} \Phi_i(t)^2\right)^2}.
\end{equation}

The IPR ranges from $1/120 \approx 0.0083$ for a perfectly delocalised
state to~$1$ for a state concentrated on a single vertex.  The
time-averaged IPR, $\langle \mathrm{IPR} \rangle_t$, is used as the
primary localisation diagnostic.

\begin{itemize}[nosep]
\item $\langle \mathrm{IPR} \rangle \lesssim 0.02$:
  effectively delocalised (comparable to uniform).
\item $0.02 \lesssim \langle \mathrm{IPR} \rangle \lesssim 0.1$:
  moderately localised (energy on $\sim$10--50~vertices).
\item $\langle \mathrm{IPR} \rangle \gtrsim 0.1$:
  strongly localised (energy on $\lesssim 10$~vertices).
\end{itemize}

% ── 5.2 ──────────────────────────────────────────────────
\subsection{Recurrence Statistics}\label{sec:recurrence}

The autocorrelation function of the field configuration is
\begin{equation}\label{eq:autocorr}
  C(\tau) = \left\langle
    \frac{\langle \Phi(t), \Phi(t+\tau) \rangle}
         {\lVert \Phi(t) \rVert \, \lVert \Phi(t+\tau) \rVert}
  \right\rangle_t.
\end{equation}

The power spectrum $S(\omega) = |\hat{C}(\omega)|^2$ reveals the
dominant frequencies and their multiplicity.  The number of distinct
spectral peaks above a threshold is denoted $N_{\mathrm{peaks}}$.

\begin{itemize}[nosep]
\item $N_{\mathrm{peaks}} = 1$: periodic.
\item $N_{\mathrm{peaks}} = 2$--5 with incommensurate ratios:
  quasi-periodic.
\item Broad-band spectrum: chaotic.
\end{itemize}

The recurrence time is defined as the smallest $\tau > 0$ for which
$C(\tau) > 1 - \delta_C$ with $\delta_C = 0.05$.

% ── 5.3 ──────────────────────────────────────────────────
\subsection{Shell Coupling}\label{sec:shell}

The 600-cell graph has palindromic distance shells
$(1, 12, 32, 42, 32, 1)$ from any vertex.  Define the shell energy
\begin{equation}\label{eq:shell_energy}
  E_d(t) = \sum_{i \in S_d(v_0)} h_i(t),
\end{equation}
where $S_d(v_0)$ is the set of vertices at graph distance~$d$ from a
reference vertex~$v_0$, and $h_i$ is the per-vertex energy.

The shell cross-correlation measures coherence between adjacent shells:
\begin{equation}\label{eq:shell_corr}
  \rho_{d,d+1}
  = \frac{\langle E_d(t) \, E_{d+1}(t) \rangle_t
    - \langle E_d \rangle \langle E_{d+1} \rangle}
    {\sigma_{E_d} \, \sigma_{E_{d+1}}}.
\end{equation}

High cross-correlation ($\rho > 0.5$) indicates coherent energy
propagation between shells.  Part~I reported $\rho \approx 0.68$ on
the 600-cell versus $\rho \approx 0.1$ on control graphs.  The
present work extends this analysis to the nonlinear regime and
examines how shell coupling varies with~$\beta$.


% ══════════════════════════════════════════════════════════
\section{Connection to Part~II Invariants}\label{sec:invariants}
% ══════════════════════════════════════════════════════════

The seven structural invariants identified in Part~II~\cite{paper2}
are properties of the static Laplacian spectrum.  The dynamical
content of these invariants---the manner in which they constrain the
attractor space---is the central concern of the present section.

\begin{enumerate}[nosep, leftmargin=*]

\item \textbf{$\vphi$-Cancellation $\to$ Balanced coupling.}
The exact cancellation of irrational spectral content
($\sum b_k = 0$, $\sum b_k d_k = 0$) ensures that the net
$\vphi$-weighted coupling across all modes vanishes.
Dynamically, this prevents systematic drift of energy toward
either $\vphi$-sector and maintains a balanced coupling
environment in which localised states can persist without
preferential leakage into irrational-frequency channels.

\item \textbf{Degree anchoring ($k=12$) $\to$ Spectral centring.}
The anchoring of irrational eigenvalues at the graph degree
($b_k < 0 \implies a_k = 12$) centres the $\vphi$-shifted
eigenvalues around the graph degree.  This constrains the
nonlinear frequency shifts~\eqref{eq:freq_shift_multi} to
remain near a common reference, facilitating mode locking.

\item \textbf{Conjugate pairing $\to$ Symmetric mode families.}
The equal multiplicities of Galois conjugate eigenvalue pairs
ensure that for every resonance channel involving a
$\vphi$-shifted frequency, a conjugate channel exists with
equal multiplicity.  This produces symmetric families of
multi-mode locked states.

\item \textbf{Square multiplicities $\to$ Backbone stability hierarchy.}
The $(\ell+1)^2$ multiplicity pattern creates a hierarchy in
which higher-multiplicity sectors ($d = 25, 36$) provide
larger basins of attraction for single-mode and multi-mode
states.  The backbone modes ($d = 1, 4, 9, 16, 25, 36$,
totalling~91) dominate the attractor landscape.

\item \textbf{Three-zone partition $\to$ Mode class separation.}
The partition $\mathcal{S} = \mathcal{S}_- \cup \mathcal{S}_0
\cup \mathcal{S}_+$ separates modes into three coupling classes.
Cross-zone nonlinear coupling is suppressed relative to
within-zone coupling, producing distinct attractor families
associated with each zone.

\item \textbf{Spectral backbone $\to$ Dominant attractors.}
The 91-mode backbone carries the dominant share of the
multiplicity-weighted spectral density.  Persistent attractors
are overwhelmingly composed of backbone modes; the 29-mode
departure contributes primarily to transient or unstable
configurations.

\item \textbf{Departure structure $\to$ Higher-energy unstable modes.}
The structured departure ($29 = 9 + 16 + 4$) provides a
finite set of modes available for energy redistribution during
transient dynamics.  These modes act as temporary energy
reservoirs rather than long-lived attractor components.

\end{enumerate}

\begin{proposition}
The spectral invariants of Part~II constrain the attractor space of
the nonlinear system.  Specifically:
\begin{itemize}[nosep]
\item The number of attractor families is bounded by the number of
  spectral sectors and their coupling structure.
\item The dominant attractors are composed of backbone modes.
\item The stability hierarchy among attractors reflects the
  multiplicity hierarchy of the spectrum.
\end{itemize}
\end{proposition}

The invariants do not uniquely determine the attractors; they constrain
the space in which attractors exist.  Paper~II provides the algebraic
constraint; the present paper identifies the dynamical realisations
within that constrained space.


% ══════════════════════════════════════════════════════════
\section{Attractor Classification}\label{sec:classification}
% ══════════════════════════════════════════════════════════

Based on the taxonomy of Section~\ref{sec:taxonomy}, the stability
analysis of Section~\ref{sec:stability}, and the diagnostic measures
of Section~\ref{sec:quantify}, four attractor classes are defined.

\begin{table}[H]
\centering
\caption{Attractor classification for the nonlinear $\Hfour$ field
system.  IPR ranges and spectral composition are time-averaged
values from the numerical simulations described in
Section~\ref{sec:numerics}.}
\label{tab:classes}
\begin{tabular}{@{}clccl@{}}
\toprule
Class & Description & IPR range &
  Spectral sectors & Stability \\
\midrule
1 & Backbone harmonic modes & $0.008$--$0.03$ &
  single ($k \leq 6$) & persistent, periodic \\
2 & Locked multi-mode states & $0.02$--$0.08$ &
  2--4 backbone sectors & persistent, quasi-periodic \\
3 & Localised breather states & $0.05$--$0.5$ &
  broad (all sectors) & persistent, quasi-periodic \\
4 & Transitional / unstable & variable &
  broad, shifting & transient \\
\bottomrule
\end{tabular}
\end{table}

\subsection{Class~1: Backbone Harmonic Modes}

These are the simplest persistent states: single-eigenvector
configurations living within the 91-mode backbone.  They are
effectively linear modes with small nonlinear frequency corrections.
Their spatial profile inherits the symmetry of the $\Hfour$ irreducible
representation to which the eigenvector belongs.

The spatial structure of Class~1 states is globally distributed across
the 120~vertices with low spatial variance.  The field amplitude
$\Phi_i$ follows a symmetry-consistent pattern determined by the
constituent eigenvector, with no strong concentration on any vertex
subset (Figure~\ref{fig:modes}, top left).

IPR is low ($\sim 1/120$ to $\sim 0.03$), placing these states near
the fully delocalised baseline.  Class~1 states dominate at low~$\beta$
($\sim$17\% of trajectories for $\beta < 0.01$) and decline to
$\sim$5--7\% in the strongly nonlinear regime
(Figure~\ref{fig:phase}), as the nonlinearity destabilises simple
harmonic configurations and redistributes energy into multi-mode
states.  Stability is high: perturbations of amplitude
$\epsilon < 0.1 \lVert \Phi^* \rVert$ return to the original state.

\subsection{Class~2: Locked Multi-Mode States}

These states involve 2--4 eigenmodes from backbone sectors with
phase-locked amplitude oscillations.  The phase locking arises from
the nonlinear frequency shifts~\eqref{eq:freq_shift_multi} and is
stabilised by the spectral compression (nine sectors rather than 120
distinct eigenvalues).

The spatial structure exhibits multi-scale interference: the
superposition of locked eigenvectors produces a pattern with partial
coherence and intermediate spatial variance
(Figure~\ref{fig:modes}, top right).  Unlike Class~1, the field
amplitude varies on multiple length scales, reflecting the
contribution of modes from different spectral sectors.

IPR is moderately elevated ($0.02$--$0.08$), reflecting partial
spatial localisation from the constructive interference of locked
modes.  Class~2 states are the most prevalent attractor family across
all values of~$\beta$, comprising $\sim$55\% of trajectories at low
nonlinearity and rising to $\sim$64\% in the strongly nonlinear regime.
This increase reflects the tendency of nonlinear frequency shifts to
bring additional mode pairs into resonance, expanding the set of
available locked configurations.  The states are quasi-periodic, with
power spectra showing 2--4 discrete peaks.

The most common locked states involve modes from sectors with the
largest multiplicities (sectors~5 and~6, with $d = 25$ and $d = 36$
respectively), consistent with the backbone stability hierarchy.

\subsection{Class~3: Localised Breather States}

These are the most distinctive configurations: spatially concentrated,
time-periodic or quasi-periodic states in which a significant fraction
of the total energy remains trapped in a small neighbourhood of the
graph.

The spatial structure is sharply localised: one or a few vertices
carry the dominant field amplitude, with a high amplitude contrast
between the localised core and the remainder of the graph
(Figure~\ref{fig:modes}, bottom left).  This localisation is not
imposed by the initial conditions but emerges dynamically from the
interplay between the nonlinear self-interaction and the constrained
spectral structure.

IPR ranges from~$0.05$ (moderate localisation) to values exceeding
$0.5$ (extreme concentration on $\sim$2--5~vertices).  At high~$\beta$,
the IPR distribution develops a heavy tail extending to $\mathrm{IPR}
> 0.3$ (Figure~\ref{fig:ipr}), well separated from the delocalised
baseline at $1/120 \approx 0.008$.  The spectral composition is
broad, spanning multiple sectors including departure modes.  However,
the time-averaged spectral content is dominated by backbone modes.

Breather states maintain a remarkably stable presence across the full
range of~$\beta$, comprising $\sim$27\% of trajectories in both the
linear and strongly nonlinear regimes.  However, the \emph{degree} of
localisation increases substantially with~$\beta$: the IPR of Class~3
states at $\beta = 5$ reaches values an order of magnitude higher than
at $\beta = 0.01$.

\subsection{Class~4: Transitional / Unstable States}

This class collects configurations that do not persist: states that
evolve toward one of the other three classes within the simulation
window, or that exhibit weakly chaotic wandering without settling into
a definite attractor.

The spatial structure is irregular, with no persistent pattern
(Figure~\ref{fig:modes}, bottom right).  The field amplitude
distribution is broad and noisy, shifting over time without converging
to a stable configuration.

Class~4 states are predominantly composed of departure-sector modes
($k = 7, 8, 9$).  Their transient nature is consistent with the
departure structure (Invariant~7 of Part~II): the 29-mode departure
provides temporary energy reservoirs but does not support long-lived
attractors.  Class~4 states are rare ($< 5\%$ of trajectories) and
peak in the crossover region between the weakly and strongly nonlinear
regimes (Figure~\ref{fig:phase}), where the competition between
linear and nonlinear organising mechanisms is most intense.

\medskip

The mapping between attractor classes and the diagnostic observables
is summarised as follows:
\begin{itemize}[nosep]
\item Class~1 $\leftrightarrow$ low IPR, dominant at low~$\beta$,
  single backbone sector.
\item Class~2 $\leftrightarrow$ intermediate IPR, phase-locked
  multi-sector composition, dominant across all~$\beta$.
\item Class~3 $\leftrightarrow$ high IPR, spatially localised,
  broad spectral content.
\item Class~4 $\leftrightarrow$ variable IPR, unstable, transient
  spectral composition.
\end{itemize}


% ══════════════════════════════════════════════════════════
\section{Results}\label{sec:numerics}
% ══════════════════════════════════════════════════════════

% ── 8.1 ──────────────────────────────────────────────────
\subsection{Overview of Numerical Experiments}\label{sec:overview}

We numerically investigate the nonlinear dynamics of the
$\Hfour$-constrained oscillator system defined on the 600-cell graph
(120~vertices, 12-regular, vertex-transitive).  The evolution is
governed by the nonlinear equation introduced in
Section~\ref{sec:system}, with particular focus on the role of the
cubic nonlinearity parameter~$\beta$.

We simulate $N = 2000$ trajectories (80~initial conditions per
$\beta$~value across 25~values of~$\beta$).  Initial conditions are
drawn from eight classes spanning localised, delocalised, and
spectrally-targeted configurations: (i)~single-vertex excitation,
(ii)~random Gaussian field ($\Phi_i \sim \mathcal{N}(0, \sigma^2)$,
$\sigma = 0.1A$), (iii)~low-mode seed (projection onto backbone
sectors $k \leq 3$), (iv)~local $\Hthree$ neighbourhood seed,
(v)~mid-mode seed (sectors 4--6), (vi)~single eigenvector with
Gaussian perturbation, (vii)~multi-vertex cluster (3--5 adjacent
vertices), and (viii)~random field with nonzero initial momentum
$\Pi_i \sim \mathcal{N}(0, \sigma_\Pi^2)$.  Initial amplitudes are
drawn uniformly from $[0.1, 2.0]$.  No energy normalisation is
applied; the total energy varies across initial conditions, reflecting
the range of dynamical regimes accessible to the system.

Time integration is performed using a velocity-Verlet (St\"ormer--Verlet)
symplectic integrator with time step $\Delta t = 0.005$.  Each
trajectory is evolved for $T = 2000$~time units ($4 \times 10^5$
steps).  The symplectic structure of the integrator ensures long-term
stability of the Hamiltonian.  Energy conservation is monitored
throughout: the relative Hamiltonian drift
$|H(t) - H(0)|/|H(0)|$ remains below $5 \times 10^{-4}$ for all
retained trajectories.  Trajectories exceeding this threshold are
discarded.

After a transient period of $T_{\mathrm{trans}} = 200$~time units,
diagnostics are computed over the remaining $1800$~time units.  The
objective is to identify dynamically stable configurations and to
characterise their structure under nonlinear evolution.

\begin{table}[H]
\centering
\caption{Parameter ranges for the numerical experiments.}
\label{tab:params}
\begin{tabular}{@{}lll@{}}
\toprule
Parameter & Symbol & Range \\
\midrule
Base frequency & $\omega_0$ & $1.0$ (fixed) \\
Laplacian coupling & $\lambda$ & $1.0$ (fixed) \\
Nonlinear strength & $\beta$ & $[10^{-3}, 5.0]$, 25 values (log-spaced) \\
Initial amplitude & $A$ & $[0.1, 2.0]$ (uniformly sampled) \\
Time step & $\Delta t$ & $0.005$ \\
Total time & $T$ & $2000$ \\
Transient period & $T_{\mathrm{trans}}$ & $200$ \\
Integration method & --- & velocity-Verlet (symplectic, 2nd order) \\
Energy drift threshold & --- & $< 5 \times 10^{-4}$ (relative) \\
Ensemble size & $N_{\mathrm{IC}}$ & $80$ per $\beta$ value (2000 total) \\
\bottomrule
\end{tabular}
\end{table}

% ── 8.2 ──────────────────────────────────────────────────
\subsection{Identification of Attractor Classes}\label{sec:id_classes}

Across all simulations, trajectories converge to a small number of
recurrent dynamical configurations.  These configurations can be
grouped into four classes:

\begin{enumerate}[nosep]
\item \textbf{Class~1:} Backbone harmonic modes.
\item \textbf{Class~2:} Locked multi-mode states.
\item \textbf{Class~3:} Localised breather-like states.
\item \textbf{Class~4:} Transitional (non-persistent) states.
\end{enumerate}

Classification is performed using quantitative criteria applied
uniformly across the ensemble:
\begin{itemize}[nosep]
\item \textbf{Class~1:} configurations with $> 80\%$ of the
  time-averaged spectral energy concentrated in a single eigenmode
  sector, and $\mathrm{IPR} < 0.03$ (near the delocalised baseline
  $1/n$).
\item \textbf{Class~2:} configurations with spectral energy
  distributed across 2--6 sectors, exhibiting stable phase
  relationships (peak autocorrelation $> 0.2$ at nonzero lag), and
  $\mathrm{IPR} < 0.10$.
\item \textbf{Class~3:} configurations with $\mathrm{IPR} > 0.05$
  and energy trapping fraction $> 0.35$ within a graph
  2-neighbourhood, indicating spatial localisation.
\item \textbf{Class~4:} configurations that fail the persistence
  criteria above over the post-transient observation window of
  $1800$~time units.
\end{itemize}

These thresholds are robust across parameter sweeps: the same four
classes emerge when thresholds are varied by $\pm 20\%$, and the
classification is stable across independent random seeds.

% ── 8.3 ──────────────────────────────────────────────────
\subsection{Representative Configurations}\label{sec:rep_configs}

Representative configurations for each attractor class are shown in
Figure~\ref{fig:modes}.

\begin{figure}[H]
\centering
\includegraphics[width=0.85\textwidth]{fig_mode_visualisation.png}
\caption{Representative field configurations from each attractor class.
Each panel shows $\Phi_i$ on the 120~vertices, ordered by distance
shell from a reference vertex.
\emph{Top left:} Class~1 (backbone harmonic).
\emph{Top right:} Class~2 (locked multi-mode).
\emph{Bottom left:} Class~3 (localised breather).
\emph{Bottom right:} Class~4 (transitional).}
\label{fig:modes}
\end{figure}

Class~1 (backbone harmonic) configurations are spatially extended,
exhibiting low variance across vertices and strong alignment with a
single Laplacian eigenmode.  These states retain the symmetry
structure of the corresponding eigenspace.

Class~2 (locked multi-mode) configurations display interference
patterns consistent with nonlinear phase-locking between multiple
eigenmodes.  These states exhibit structured spatial variation with
partial coherence across the graph.

Class~3 (breather-like) configurations are spatially localised, with
a significant fraction of the total amplitude concentrated on a small
subset of vertices.  These configurations persist over long times and
are consistent with intrinsic localised modes observed in nonlinear
lattice systems~\cite{flach2008,aubry1997}.

Class~4 (transitional) configurations are irregular and do not exhibit
temporal persistence or coherent structure.  These trajectories
typically evolve toward one of the other classes or decay into
fluctuating states.

Importantly, these configurations arise dynamically from the system
evolution and are not imposed by construction.

% ── 8.4 ──────────────────────────────────────────────────
\subsection{Dependence on Nonlinearity}\label{sec:beta_dep}

The distribution of attractor classes as a function of the nonlinear
parameter~$\beta$ is shown in Figure~\ref{fig:phase}.

\begin{figure}[H]
\centering
\includegraphics[width=0.85\textwidth]{fig_phase_diagram.png}
\caption{Attractor class distribution as a function of the nonlinear
strength~$\beta$ (log scale), from an ensemble of 80~trajectories per
$\beta$~value across 25~values spanning $[10^{-3}, 5]$.  Three regimes
are indicated by dashed lines.}
\label{fig:phase}
\end{figure}

In the \emph{linear regime} ($\beta \ll 1$), dynamics are dominated
by individual eigenmodes, and Class~1 configurations are prevalent.
Mode coupling is minimal, and trajectories remain close to linear
superpositions.

In the \emph{weakly nonlinear regime}, multi-mode interactions become
significant.  Class~2 configurations increase in frequency, reflecting
the emergence of phase-locked combinations of eigenmodes.

In the \emph{strongly nonlinear regime} ($\beta \gtrsim 1$), localised
configurations (Class~3) become more prominent.  At the same time, the
fraction of purely harmonic configurations decreases.

The distribution of attractor classes varies continuously with~$\beta$,
indicating a smooth transition between dynamical regimes rather than a
sharp bifurcation.

% ── 8.5 ──────────────────────────────────────────────────
\subsection{Localisation Quantification via IPR}\label{sec:ipr_results}

To quantify localisation, we compute the inverse participation ratio
(IPR) as defined in equation~\eqref{eq:ipr}:
\[
  \mathrm{IPR} = \frac{\sum_i |\Phi_i|^4}
    {\left(\sum_i |\Phi_i|^2\right)^2}.
\]
For a fully delocalised configuration,
$\mathrm{IPR} = 1/n = 1/120$.

The distribution of time-averaged IPR values across nonlinear regimes
is shown in Figure~\ref{fig:ipr}.

\begin{figure}[H]
\centering
\includegraphics[width=0.85\textwidth]{fig_ipr_distribution.png}
\caption{Distribution of time-averaged IPR across all trajectories for
three representative values of~$\beta$.  The dotted line marks the
fully delocalised baseline $\mathrm{IPR} = 1/n = 1/120$.}
\label{fig:ipr}
\end{figure}

At low~$\beta$, the IPR distribution is tightly concentrated near the
baseline value, consistent with spatially extended modes.  At
intermediate~$\beta$, the distribution broadens, indicating increased
spatial variation and partial localisation.  At high~$\beta$, a
pronounced heavy tail appears at large IPR values, indicating the
emergence of strongly localised configurations.

The emergence of a high-IPR tail at larger~$\beta$ indicates the
formation of spatially localised dynamical states.  This behaviour is
consistent with known properties of nonlinear lattice systems
supporting intrinsic localised modes (discrete
breathers)~\cite{flach2008,campbell2004}.

% ── 8.6 ──────────────────────────────────────────────────
\subsection{Summary of Results}\label{sec:results_summary}

The $\Hfour$-constrained nonlinear system exhibits a small number of
persistent dynamical configurations that can be classified into four
attractor classes.  The relative prevalence of these classes depends
smoothly on the nonlinear coupling strength~$\beta$.  In particular,
increasing nonlinearity leads to a transition from globally distributed
harmonic modes to spatially localised configurations.

These results indicate that the system possesses a small and structured
set of recurrent attractor classes within the explored parameter range.

% ── 8.7 ──────────────────────────────────────────────────
\subsection{Comparison with Control Graphs}\label{sec:baseline}

We compare the $\Hfour$ graph against an ensemble of random 12-regular
graphs matched in size ($n = 120$) and degree ($k = 12$), using
identical dynamical parameters, initial-condition sampling, integration
settings, and classification criteria.  The purpose of this comparison
is to determine whether the observed attractor structure is a generic
consequence of nonlinear dynamics on regular graphs, or whether it
depends on the specific spectral organisation of the $\Hfour$ graph.

The same classification criteria, thresholds, and feature extraction
pipeline were applied to both the $\Hfour$ graph and all control graphs
without graph-specific tuning.  Here ``structured'' refers to
configurations exhibiting temporal persistence and nontrivial spatial
organisation (including multi-mode locking or localisation), as opposed
to transient or irregular dynamics.

The random regular graphs (RRGs) are generated by configuration-model
pairing with local repair and verified to be simple, connected,
12-regular graphs.  Results for the control ensemble are reported as
mean $\pm$ standard deviation over $M = 10$ independent graph
realisations.  Each RRG has 120~distinct eigenvalues (no degeneracy),
compared with 9 for the 600-cell
(Figure~\ref{fig:baseline_spectrum}).

\begin{figure}[H]
\centering
\includegraphics[width=0.75\textwidth]{fig_baseline_spectrum.png}
\caption{Number of distinct Laplacian eigenvalues for the $\Hfour$
graph (index~0, red) and 10~random 12-regular graphs (blue).  The
600-cell has 9~distinct eigenvalues; all control graphs have~120.}
\label{fig:baseline_spectrum}
\end{figure}

The $\Hfour$ graph exhibits a substantially higher fraction of
persistent structured configurations, while trajectories on the
control graphs are predominantly classified as transitional under the
same criteria (Figure~\ref{fig:baseline_classes}).  On the $\Hfour$
graph, Classes~1 and~2 account for $\sim$70\% and Class~3 for
$\sim$27\%.  On the RRGs, $\sim$96\% of trajectories are classified as
Class~4 (transitional), with Classes~1--3 comprising fewer than 4\%
combined.

\begin{figure}[H]
\centering
\includegraphics[width=0.95\textwidth]{fig_baseline_classes.png}
\caption{Attractor class distribution as a function of~$\beta$ for
the $\Hfour$ graph (circles, solid) and random 12-regular graphs
(squares, dashed; shaded region shows $\pm 1\sigma$ across $M = 10$
controls).  On the $\Hfour$ graph, trajectories distribute across
all four classes.  On control graphs, nearly all trajectories are
classified as Class~4 (transitional).}
\label{fig:baseline_classes}
\end{figure}

\begin{figure}[H]
\centering
\includegraphics[width=0.95\textwidth]{fig_baseline_ipr.png}
\caption{IPR distributions for the $\Hfour$ graph (red) and pooled
random regular graphs (blue) at three values of~$\beta$.  The
$\Hfour$ graph shows broader IPR distributions with a tail extending
to higher values, indicating stronger spatial localisation.}
\label{fig:baseline_ipr}
\end{figure}

\noindent This contrast is reflected not only in classification
outcomes but also in graph-invariant observables:
\begin{itemize}[nosep]
\item \textbf{Spectral compression:}
  $\Hfour$: 9~distinct eigenvalues; RRG: $120 \pm 0$ (no degeneracy).
\item \textbf{Structured attractors:}
  $\Hfour$: $\sim$98\% of trajectories in Classes~1--3; this fraction
  was not observed in the tested control ensemble ($< 4\%$).
\item \textbf{Localisation:}
  mean IPR on $\Hfour$ ($0.062$) exceeds the RRG value ($0.028$) by
  a factor of $\sim$2.3, with the $\Hfour$ distribution exhibiting a
  substantially heavier tail at high~$\beta$.
\item \textbf{Attractor diversity:}
  the Shannon entropy of the class distribution (averaged over~$\beta$)
  is $\sim$1.5~bits for the $\Hfour$ graph versus $\sim$0.2~bits for
  the RRG ensemble, indicating that the $\Hfour$ graph supports a
  richer distribution of distinct attractor types rather than
  collapsing into a single dynamical class.
\end{itemize}

Because the primary classification includes spectral components
(eigenmode projections), care must be taken to ensure that the observed
differences are not solely driven by the compressed eigenspace
structure of the $\Hfour$ graph.  To address this, we repeat the
classification using only graph-invariant observables --- IPR,
temporal persistence (autocorrelation), and energy variance --- without
reference to the eigenbasis.

Under graph-invariant classification, the $\Hfour$ graph retains
$\sim$99\% of trajectories in structured classes (C1--C3), with
$\sim$44\% classified as spatially localised (Class~3).  The control
RRGs retain $\sim$82\% in structured classes but produce localised
trajectories at an approximately two orders of magnitude lower rate
($\sim$0.3\% Class~3).

In particular, a degree-preserving rewiring of the 600-cell
graph---destroying its $\Hfour$ symmetry while maintaining identical
size and degree---eliminates the structured attractor distribution and
restores behaviour comparable to random regular graphs ($\sim$0.3\%
Class~3, $\sim$16\% Class~4, 120~distinct eigenvalues).  This isolates
the $\Hfour$ symmetry, rather than the degree sequence or graph size,
as the relevant structural factor.

Taken together, these results provide consistent evidence that the
observed attractor structure is not an artefact of the classification
procedure and is not reproduced in degree-matched control graphs under
identical conditions.

The consistency of these observations across spectral, dynamical, and
graph-invariant analyses supports the conclusion that the attractor
structure is tied to the specific symmetry and spectral properties of
the $\Hfour$ graph.  These results are presented as properties of a
nonlinear dynamical system on a structured graph and do not assume any
direct mapping to physical observables.


% ══════════════════════════════════════════════════════════
\section{Discussion}\label{sec:discussion}
% ══════════════════════════════════════════════════════════

% ── 9.1 ──────────────────────────────────────────────────
\subsection{Relation to Spectral Structure}\label{sec:disc_spectral}

The observed dynamical behaviour reflects the highly constrained
spectral structure of the $\Hfour$ graph Laplacian described in
Part~II of this series~\cite{paper2}.

The strong degeneracy of eigenvalues organises the linear modes into a
small number of sectors.  The dominance of Class~1 and Class~2
configurations is consistent with the existence of a spectral backbone
(the 91-mode subspace spanning sectors~1--6), while the emergence of
Class~3 configurations reflects nonlinear interactions across these
sectors.

The finite number of distinct eigenvalue sectors appears to limit the
number of dynamically stable configurations.

% ── 9.2 ──────────────────────────────────────────────────
\subsection{Interpretation as a Constrained Nonlinear System}
\label{sec:disc_constrained}

The system studied here is a nonlinear oscillator network defined on
a highly symmetric graph.  The combination of spectral compression and
nonlinear coupling leads to a restricted set of persistent dynamical
states, suggesting a strong influence of the underlying spectral
structure on the attractor taxonomy.

The baseline comparison of Section~\ref{sec:baseline} supports this
interpretation: on random 12-regular graphs with identical size,
degree, and dynamics, $\sim$96\% of trajectories are classified as
Class~4 (transitional), compared with $\sim$2\% on the $\Hfour$ graph.
This contrast was not observed in any of the tested control graphs,
suggesting that the structured attractor space is non-generic within
the tested ensemble of 12-regular graphs and reflects the spectral
organisation of the $\Hfour$ graph.

% ── 9.3 ──────────────────────────────────────────────────
\subsection{Localisation Mechanism}\label{sec:disc_localisation}

The emergence of localised configurations can be understood in terms
of nonlinear frequency shifts and energy trapping.

The cubic nonlinearity introduces amplitude-dependent frequency
corrections, enabling resonance conditions between
modes~\cite{flach2008,mackay1994}.  Combined with the finite
connectivity of the graph ($k = 12$), this allows energy to become
localised in space.

Localised configurations arise without disorder, consistent with the
theory of discrete breathers in nonlinear lattice
systems~\cite{aubry1997,mackay1994}.

% ── 9.4 ──────────────────────────────────────────────────
\subsection{Structured Attractor Space}\label{sec:disc_finite}

Within the explored parameter range, the system exhibits a small and
structured set of recurrent attractor classes.

Stability is assessed numerically through persistence over long time
windows ($1800$~time units post-transient) and robustness to small
perturbations in initial conditions, though a formal Lyapunov analysis
is not performed in this work.

While this classification is based on numerical evidence and is not
proven to be exhaustive, it is stable across initial conditions and
parameter sweeps.  This suggests that the attractor space of the system
is strongly constrained by the underlying symmetry and spectral
structure.

% ── 9.5 ──────────────────────────────────────────────────
\subsection{Limitations}\label{sec:limitations}

The present results are numerical and do not constitute a complete
analytical classification of the system.

Limitations include:
\begin{itemize}[nosep]
\item finite simulation time and sampling;
\item dependence on chosen parameter ranges;
\item absence of a formal stability proof for each attractor class;
\item lack of an analytical derivation of the attractor structure;
\item the control study is limited to a finite ensemble of random
  12-regular graphs and does not exclude the possibility that other
  highly structured non-$\Hfour$ graphs may exhibit related behaviour.
\end{itemize}

No direct mapping to physical observables is assumed in this work.

% ── 9.6 ──────────────────────────────────────────────────
\subsection{Outlook}\label{sec:outlook}

Future work should focus on:
\begin{itemize}[nosep]
\item analytical derivation of attractor classes using representation
  theory of~$\Hfour$;
\item stability analysis via perturbative and variational methods;
\item extension to coupled $\Hfour$ systems and higher-dimensional
  embeddings;
\item systematic exploration of parameter space.
\end{itemize}

\bigskip

\noindent\emph{The results indicate that a nonlinear oscillator system
defined on the $\Hfour$ symmetry graph exhibits a structured and
reproducible set of dynamical attractors, whose organisation reflects
the underlying spectral constraints of the graph.}


% ══════════════════════════════════════════════════════════
\section{Conclusion}\label{sec:conclusion}
% ══════════════════════════════════════════════════════════

The numerical experiments reported here indicate that the nonlinear
field system defined by~\eqref{eq:eom} on the 600-cell vertex graph
admits a small and structured set of recurrent dynamical
configurations within the explored parameter range.  Four attractor
classes have been identified and characterised: backbone harmonic
modes, phase-locked multi-mode states, spatially localised
breather-like states, and transitional non-persistent states.

The attractor taxonomy is consistent with the seven structural
invariants of the $\Hfour$ Laplacian spectrum identified in Part~II.
The spectral compression (nine sectors across 120~degrees of freedom),
the backbone--departure decomposition (91~+~29), the $\vphi$-zone
partition, and the conjugate pairing structure collectively constrain
the space of persistent nonlinear configurations to the classes
described here.

The relationship between the three parts of this series may be
summarised as follows:
\begin{center}
\begin{tabular}{@{}ll@{}}
Part~I & establishes the phenomenon (coherent dynamics on $\Hfour$ geometry) \\
Part~II & establishes the structure (spectral invariants governing the spectrum) \\
Part~III & establishes the taxonomy (attractor families realised within the \\
  & constrained spectral space) \\
\end{tabular}
\end{center}

Together, the three papers indicate that the 600-cell graph possesses
a non-generic spectral algebra that produces structured nonlinear
dynamics with a classifiable set of recurrent attractor families.  No
physical interpretation is required for or assumed in these
conclusions.

\paragraph{Reproducibility.}
All simulations use the velocity-Verlet symplectic integrator applied
to the explicit 600-cell graph Laplacian.  The simulation pipeline,
parameter ranges, and diagnostic code are specified in full in the
accompanying repository.  No external datasets or proprietary
software are required.


% ══════════════════════════════════════════════════════════
% References
% ══════════════════════════════════════════════════════════
\begin{thebibliography}{10}

\bibitem{paper1}
L.~Smart,
``Symmetry-Constrained Resonant Modes on $H_4$ Geometry:
A Geometric Framework for Vibrational Field Dynamics,''
\textit{VFD Working Paper Series}, WO-VFD-PAPER-001, 2025.

\bibitem{paper2}
L.~Smart,
``Algebraic Invariants of the $H_4$ Laplacian Spectrum,''
\textit{VFD Working Paper Series}, WO-VFD-PAPER-002, 2025.

\bibitem{flach2008}
S.~Flach and A.~V.~Gorbach,
``Discrete breathers --- Advances in theory and applications,''
\textit{Physics Reports} \textbf{467}, 1--116, 2008.

\bibitem{aubry1997}
S.~Aubry,
``Breathers in nonlinear lattices: Existence, linear stability
and quantization,''
\textit{Physica D} \textbf{103}, 201--250, 1997.

\bibitem{mackay1994}
R.~S.~MacKay and S.~Aubry,
``Proof of existence of breathers for time-reversible or
Hamiltonian networks of weakly coupled oscillators,''
\textit{Nonlinearity} \textbf{7}, 1623--1643, 1994.

\bibitem{campbell2004}
D.~K.~Campbell, S.~Flach, and Y.~S.~Kivshar,
``Localizing energy through nonlinearity and discreteness,''
\textit{Physics Today} \textbf{57}, 43--49, 2004.

\bibitem{kevrekidis2009}
P.~G.~Kevrekidis,
\textit{The Discrete Nonlinear Schr\"odinger Equation:
Mathematical Analysis, Numerical Computations and Physical
Perspectives}.
Springer, 2009.

\bibitem{chung1997}
F.~R.~K.~Chung,
\textit{Spectral Graph Theory}.
AMS, 1997.

\bibitem{coxeter1973}
H.~S.~M.~Coxeter,
\textit{Regular Polytopes}, 3rd~ed.
Dover Publications, 1973.

\bibitem{humphreys1990}
J.~E.~Humphreys,
\textit{Reflection Groups and Coxeter Groups}.
Cambridge University Press, 1990.

\end{thebibliography}

\vfill
\noindent\rule{\textwidth}{0.4pt}\\[4pt]
{\footnotesize
\textit{Vibrational Field Dynamics Working Paper Series}
\hfill Part~III \quad v1.7 (final-hardened)}

\end{document}
